% Article-ready mathematics for the conditional-G composition.
% Requires amsmath, amssymb, amsthm; theorem environment names can be adjusted.
\section{A logical composition theorem for the conditional G construction}
\label{sec:conditional-g}

The new combination uses two distinct existing improvements. The conditional
K source changes the outer set-theoretic schema. The G source merges equality
propagation with a restricted equality-vacuity helper. Their compatibility is
not a consequence of subtracting their state savings: it follows first from a
uniform deductive-equivalence theorem, and then from a separately checked
source compiler and transition table.

Write $I(A,B)=A\mathbin{\to}B$, $Q_iA=\operatorname{All}(i,A)$,
$d_{ab}=\operatorname{Eq}(a,b)$ and $r_{ab}=\operatorname{Rel}(a,b)$.
The designated falsehood is the exact native code $F=1=r_{00}$.
Put $N(A)=I(A,F)$ and $\exists_iA=N(Q_iN(A))$. A closed proof means a
proof without undischarged assumptions; its formula may contain free labels.
Generalization is applied only to such closed proofs.

\subsection{Uniformity in the core and in the helper binder}

Let $\kappa\colon\mathbb N^2\to\mathbb N$ be any total function for the
moment, and define
\[
 H_\kappa(a,b)=d_{ab}\to Q_{\kappa(a,b)}d_{ab},\qquad
 G_\kappa(a,b,c)=d_{ca}\to
       \bigl(d_{cb}\to Q_{\kappa(a,b)}d_{ab}\bigr).
\]
The eventual machine uses
\[
 \kappa(a,b)=4P(a,b)+1,\qquad
 P(a,b)=\frac{(a+b)(a+b+1)}2+b.
\]
This binder exceeds both endpoint labels, but the syntactic fusion theorem
does not require that inequality. Freshness is needed when calibrating the
helper to ordinary equality semantics, not when proving its equivalence to
the merged G row.

Let $C(a,b,c)$ be an arbitrary family of native formula codes. It may be an
unconditional K core, a conditional K core, or any other expression. Let
$V(a,b,c)$ be another arbitrary family. For the effective version, assume
that these native-code functions, $\kappa$, and the gate witness are primitive
recursive, as they are in the present construction. Without this effectiveness
assumption the same argument still gives deductive equivalence.
Consider two numerical calculi with
the same truth seed, modus ponens, closed generalization, and retained rows
$\rho_1,\ldots,\rho_6,\rho_8$. The old calculus additionally has
\[
 \rho_7(a,b,c)=d_{ca}\to(d_{cb}\to d_{ab})
\]
for every triple and, whenever $c\notin\{a,b\}$, the bundle
\[
 B_{\mathrm{old}}(a,b,c)
   =N\bigl(I(C(a,b,c),I(V(a,b,c),N(H_\kappa(a,b))))\bigr).
\]
The new calculus replaces $\rho_7$ by the \emph{ungated} row $G_\kappa$
and replaces each enabled bundle by
\[
 B_{\mathrm{new}}(a,b,c)
   =N\bigl(I(C(a,b,c),N(V(a,b,c)))\bigr).
\]
Disabled bundles remain the native truth seed $0$ in both calculi.

\paragraph{Theorem (core-independent equality fusion).}
The two calculi derive exactly the same native formula codes. The proof
transformations are uniform finite syntactic transformations. In particular,
their distinguished refutation predicates, derivability of code $1$, agree.
The assertion remains valid for any common gate with an effective witness
$w(a,b)$ satisfying the gate for every ordered pair; here one may take
$w(a,b)=\max(a,b)+1$.

\paragraph{Proof.}
The unchanged $\rho_8/\rho_6$ bridge supplies classical propositional logic
with designated falsehood $F$. In detail, the retained condensed-detachment
derivations give K, S and Peirce from $\rho_8$; $\rho_6$ and K give
$F\to A$; classical negation and double-negation elimination follow. The
argument does not assume equality reflexivity, equality vacuity, G, or any
set-theory axiom.

Closed generalization and $\rho_2$ give universal monotonicity, and
contraposition gives existential monotonicity. One also has, for every
label $i$ and arbitrary native formula code $A$,
\begin{equation}
 Q_iA\longrightarrow\exists_iA.
 \tag{S}
\end{equation}
Indeed $\rho_1(i,i)$ gives $\exists_i d_{ii}$. Generalizing
$A\to(d_{ii}\to A)$ and distributing gives
$Q_iA\to Q_i(d_{ii}\to A)$. The dual consequence of $\rho_2$ is
$Q_i(d_{ii}\to A)\to(\exists_i d_{ii}\to\exists_iA)$, so propositional
composition proves (S). Consequently
\begin{equation}
 Q_iQ_iA\longrightarrow\exists_iQ_iA\longrightarrow A,
 \tag{D}
\end{equation}
using $\rho_6(i,A)$ in the last step. We have not assumed universal
elimination $Q_iA\to A$ for arbitrary codes.

Fix $x$ and set $z=\kappa(x,x)$ and $e=d_{xx}$. The G instance
$G_\kappa(x,x,z)$ contracts propositionally to $d_{zx}\to Q_ze$.
Existential monotonicity and $\rho_1(z,x)$ yield $\exists_zQ_ze$, so
$\rho_6(z,e)$ proves $d_{xx}$. This proves reflexivity without using the
deleted helper or $\rho_7$. Now $G_\kappa(a,b,a)$ and the proved
$d_{aa}$ yield the exact helper $H_\kappa(a,b)$.

Put $i=\kappa(a,b)$ and $e=d_{ab}$. Generalize the closed helper and
distribute to obtain $Q_ie\to Q_iQ_ie$. Composing with (D) gives
$Q_ie\to e$. Applied to the consequent of $G_\kappa(a,b,c)$, this
recovers the literal old $\rho_7(a,b,c)$. Conversely old $\rho_7$
followed by old $H_\kappa$ immediately proves G.

The exact negative bundles satisfy the propositional implications
\[
 B_{\mathrm{old}}\to B_{\mathrm{new}},\quad
 B_{\mathrm{old}}\to H_\kappa,\quad
 B_{\mathrm{new}}\to(H_\kappa\to B_{\mathrm{old}}).
\]
For any $a,b$, choose $c=\max(a,b)+1$ to extract the needed helper from
an enabled old bundle. This extraction is independent of the truth or form
of $C(a,b,c)$: a bundle axiom is a conjunction, so its helper conjunct is
available even when the conditional core's seed premise is false. In the
new calculus, G supplies the helper directly. Replace changed axiom leaves
by these finite closed derivations and preserve MP and closed GEN. Induction
on a proof tree gives both proof transformations.\hfill$\square$

\subsection{The exact conditional source formula}

The source compiler works with the fused constructor
\[
 X(i,A,B)=N\bigl(Q_i I(A,N(B))\bigr),
\]
which is provably equivalent in the retained logical basis to
$\exists_i(A\wedge B)$. At native parameters $a,b,c$ put $m=c$ and
$k=c$. The \emph{same number} is the opaque matrix code and the private
container-variable label. Define the following constructor expressions:
\begin{align*}
 r&=r_{bk},&
 S&=Q_k(r_{ka}\to r_{kb}),\\
 W&=Q_b(m\to X(b,m,r)),&
 U&=Q_a(S\to W),\\
 H&=Q_b(r\to X(a,m,U)),&
 C_c&=m\to X(k,r,H).
\end{align*}
Full membership vacuity remains
\[
 V_R(a,b,c)=r_{ab}\to Q_c r_{ab}.
\]
The new source emits the eight old logical row positions, with G in the
$\rho_7$ position, and the gated two-component bundle
\begin{equation}
 B_{c,G}(a,b,c)=
 \begin{cases}
  0,& c=a\text{ or }c=b,\\
  N\bigl(I(C_c,N(V_R(a,b,c)))\bigr),&\text{otherwise}.
 \end{cases}
 \tag{CG}
\end{equation}
Its pre-fusion parent has the same conditional core and membership vacuity,
old $\rho_7$, and the additional $H_\kappa$ conjunct. Applying the theorem
with $C=C_c$ and $V=V_R$ proves deductive equivalence with that precise
conditional-223 source theory. There is no return to the unconditional core
at this step, and no invalid assumption that conditional and unconditional
K have identical theorems.

The disabled gate is unchanged. An enabled bundle in (CG) is an implication
into $F=1$, hence its native code is positive and divisible by four; it can
be neither $0$ nor $1$. The G row is also an implication with a positive
quantified consequent. Thus the inherited non-target facts used by the
proof enumerator remain applicable. The number and positions of logical
rows, stack pushes, and fields consumed by AX remain unchanged.

\subsection{Ordinary conditional K and its consistency calibration}

For clarity, the conditional core corresponds to the ordinary schema
\[
 \forall a\,\forall x\,
   \bigl(\varphi(x,a,\vec q)\to
          \exists C\,(x\in C\wedge H_\varphi(C))\bigr),
\]
where
\begin{align*}
 W_\varphi(p,C)&=
   (\exists y\,\varphi(y,p,\vec q))\to
     \exists y\,(y\in C\wedge\varphi(y,p,\vec q)),\\
 H_\varphi(C)&=\forall u\,
 \left(u\in C\to
 \left((\exists p\,\varphi(u,p,\vec q))\wedge
       \forall p\,(p\subseteq u\to W_\varphi(p,C))\right)\right).
\end{align*}
Let $K_c$ be this schema plus genuine equality; add irreflexivity
$\mathrm{IR}:\forall x\,x\notin x$. The inherited conditional-K argument
establishes a uniform interpretation of the unconditional K schema plus IR
inside $K_c+\mathrm{IR}$. The essential steps are recalled because
conditional seeding alone does not prove the existence of a raw empty object.

Conditional K first constructs exact singletons and pairs. A singleton of
$a$ comes from the matrix $y=a$. For a pair $\{a,b\}$ choose a singleton
$s=\{a\}$, so IR implies $s\ne a$, and use
$(p=s\wedge y=a)\vee(p=a\wedge y=b)$ with seed $a$ at input $s$.
Closure at the eligible input $p=a\subseteq a$ forces $b$ into the
container, and support excludes all other members.

Choose distinct markers $t,s$ with $s=\{t\}$. Put
$Q(x)\equiv(t\in x\wedge s\in x)$ and let
\[
 \operatorname{Adm}(T)\equiv
 \forall u\in T\,
 \left(Q(u)\wedge\forall v\in u\,
       (v=t\vee v=s\vee v\in T)\right).
\]
The interpreted domain is the finite, nonrecursive formula
$D(x)\equiv\exists T\,(x\in T\wedge\operatorname{Adm}(T))$.
It excludes both markers, is hereditary away from them, and contains the
exact pair $\{t,s\}$, which is empty in the interpreted domain.

The key closure lemma says that if $Q(C)$ holds and every non-marker
member of $C$ belongs to $D$, then $D(C)$. Apply conditional K to
$\operatorname{Sing}(y,p)\wedge(D(y)\vee y=C)$ with a singleton input
for $C$. Support gives $Q$ on every resulting member; closure at singleton
inputs puts every needed non-marker member into one common witnessing set.
This avoids a hidden collection or choice of individual $D$-witnesses.

To transfer an arbitrary $D$-relativized instance, choose singleton inputs
$p_t=\{t\}$ and $p_s=\{s\}$, both outside $D$, and use the raw matrix
\[
 (p=p_t\wedge y=t)\vee(p=p_s\wedge y=s)
 \vee\bigl(D(p)\wedge D(y)\wedge\varphi^D(y,p,\vec q)\bigr).
\]
Starting from an interpreted output forces both markers into the new
container. Support puts every other member in $D$, and the closure lemma
puts the container itself in $D$. An interpreted subset is a raw subset:
its marker members are already in every interpreted containing object,
and all other members are hereditary in $D$. Thus the interpreted domain
satisfies conditional K and has an empty object. Splitting on the seed
premise now proves unconditional K there. The reverse direction, from
unconditional K to conditional K, simply retains the old container when
the seed premise holds. These are effective finite-schema interpretations,
so the inherited K/ZFC calibration yields
\[
 \operatorname{Con}(K_c+\mathrm{IR})
 \longleftrightarrow\operatorname{Con}(\mathrm{ZFC}).
\]

\subsection{The numerical bridge and its freshness obligations}

The restricted-helper calibration remains essential: deleting its conjunct
is justified here only because G derives the exact helper. The inherited
proof interprets primitive equality by membership indiscernibility while
preserving membership, implication, and every binder. It supplies all old
logical rows and their auxiliary-label transport, including coincidences.

For the lower direction choose an ordinary matrix $f$ with input label $1$
and output label $3$, apply the equality interpretation to obtain
$\theta(f)$, double-negation-pad it, and \emph{then} choose
\[
 c=\operatorname{enc}(\neg\neg\theta(f)).
\]
The constructor-code induction gives $c>3$ and makes $c$ larger than
every label in the padded interpreted matrix. Choose the comparison
ordinary instance's private variable to be this $c$. The K skeleton is
equality-free and commutes literally with $\theta$, so its interpreted
matrix has actual native code $c$ and is exactly the source instance with
$k=m=c$. This order of choices matters: choosing $c$ from $f$ before
interpreting equality would generally give the wrong private name. The
translated IR becomes the generalized identity $Q_0(F\to F)$.

For the upper direction the inherited decoder theorem proves private-label
freshness for every matrix code, including $0$ and $1$. One must not replace
it by the false assertion that every label in the decoded formula for code
$0$ is strictly below $0$. At enabled triples the decoded unconditional
core entails the conditional core by the fresh-container quantifier rule;
the inherited coincident-role and distinct-role cases are retained. G may
either be checked in genuine equality semantics using the freshness of
$\kappa(a,b)$, or, more directly, transported through the literal
deductive-equivalence theorem back to the calibrated conditional source.

All these transformations concern finite formulas and finite proof trees.
Fresh names are selected by explicit maxima and the source encoding is a
primitive-recursive function. They are formalizable in ZFC and yield its
internal equivalence between the two finite-refutation predicates. They do
not prove consistency of ZFC or unconditional non-halting of the machine.

\subsection{Schematic certificate and remaining proof boundary}

The new certificate \texttt{uniform\_g\_proof.json} has 57 proof nodes,
53 formula expressions, and eight requested endpoints. Unlike a list of
concrete numerical examples, its labels are formal terms generated by
$a,b,c$ and an uninterpreted binary function $\kappa$. Its arbitrary core
is a formula placeholder. A structural substitution induction shows that
every evaluation of the label terms in natural numbers, including all
coincidences and zero, maps each checked node to a permitted native proof
node. The same argument permits replacement of the core placeholder by the
exact conditional core above.

The generator/checker and the separate checker
\texttt{verify\_uniform\_g.py} validate all MP and closed-GEN edges, exact
schema leaves, exact root endpoints, and the direction-specific absence of
circular old/G premises. Propositional leaves are independently truth-table
checked; they abbreviate the inherited $\rho_8/\rho_6$ classical bridge.
Negative controls reject a circular helper premise, a changed G instance,
and substitution of a different proved root for a requested endpoint.

This certificate strengthens the executable evidence for the new logical
composition. The ordinary two-marker interpretation, the restricted-helper
and numerical K calibration, exact source binding, complete blank startup,
and universal binary runtime are separate links. The result is not a new
compiled Lean or Rocq endpoint. A literal machine theorem must compose those
links for the exact delivered transition table.
